\section{St. Venant Torsional Warping}\label{sec:torsion-warping}

The objective of this section is to detail the computation of St. Venant's warping function \(\varphi\) and related cross section parameters \(C_{\varphi}, C_w, J, ...\).

\subsection{The Neumann Problem}
Consider a body with reference configuration 
\begin{equation*}
    \bm{x}_0(\reflenx) + \SectionCoordinate(\varsigma)
\end{equation*}
undergoing a deformation with the form
\begin{equation*}
    \bm{x}(\reflenx, \varsigma ) = \bm{x}_0(\reflenx) + \SectionCoordinate(\varsigma) + \vartheta \, \FrameBasis\times(\SectionCoordinate - \bm{\SectionLocation}) + \alpha \varphi(\varsigma) \FrameBasis
\end{equation*}
Subject to the traction-free boundary condition
\begin{equation*}
\begin{aligned}
    \bm{S}\mathbf{n} &= \bm{0}
\end{aligned}
\end{equation*}
where \(\mathbf{n}\) is the unit normal on the boundary in the section's plane. 
The displacement field \(\bm{u}\) is 
\begin{equation*}
\bm{u}(\reflenx, y, z) =
\vartheta \, \FrameBasis\times(\SectionCoordinate - \bm{\SectionLocation}) + \alpha \varphi(\SectionCoordinate) \FrameBasis
\end{equation*}
and the infinitesimal strain vector
\begin{equation*}
\bm{\varepsilon} = \vartheta' \FrameBasis \times (\SectionCoordinate - \bm{\SectionLocation}) + \alpha' \varphi \FrameBasis + \alpha \nabla \varphi
\end{equation*}
The traction-free boundary condition implies
\begin{equation*}
    \nabla_{\mathbf{n}}\bm{\varepsilon} = \bm{0}
\end{equation*}
Recall \cref{eq:saint-venant-stress} and that for a prismatic section the outward normal is orthogonal to the axial basis \(\FrameBasis\), \(\mathbf{n}\cdot \FrameBasis = 0\).
\begin{marsdenbox}
\noindent
Torsional warping problem for rotation about some point \(\bm{\SectionLocation}\)
\begin{equation}
\begin{array}{rll}
\Delta \varphi_{\SectionLocation} &=0 & \text { in } \Omega, \\
\nabla_{\mathbf{n}} \varphi_{\SectionLocation}  &=-\alpha \FrameBasis \times (\mathbf{r}-\bm{\SectionLocation}) \cdot \mathbf{n} & \text { on } \partial \Omega  \\
\SectionIntegral{\varphi_{\SectionLocation}} &= 0
\end{array}
\label{eq:torsion-neumann}
\end{equation}
where $\mathbf{n}$ is the unit outward normal to the boundary of the cross section domain.
\end{marsdenbox}


\subsection{Center of Twist}
A natural choice for \(\bm{\SectionLocation}\) is to select it such that the following conditions hold:
\begin{equation}
    \SectionIntegral{\varphi_{\SectionLocation} (\bm{r}-\bm{\SectionLocation})\times\FrameBasis} = \bm{0}
    \label{eq:shear-center-constraint}
\end{equation}
Thus, the point \(\tilde{\bm{\SectionLocation}}\) is defined as that for which the corresponding solution \(\tilde{\varphi}\) to \cref{eq:torsion-neumann} with \(\pi = \tilde{\bm{\SectionLocation}}\) satisfies \cref{eq:shear-center-constraint}. 
Note that \cref{eq:shear-center-constraint} simply requires that the \(m-w\) coupling tensor of \cref{eq:def-linear-isotropic-coupling-tensors} vanishes.
...
say
\[
\tilde{\varphi} \triangleq \varphi - [\tilde{\bm{\SectionLocation}} , \FrameBasis, (\bm{r} - \hat{\bm{\SectionLocation}})]
\]
In components 
\[
\tilde{\varphi} = \varphi - 
\]
or equivalently
\[
\tilde{\varphi}\bm{r}\times\FrameBasis = \varphi\bm{r}\times\FrameBasis - \bm{r}\times\FrameBasis \, \left(\tilde{\bm{\SectionLocation}} \cdot \FrameBasis \times \hat{\bm{r}}\right)
\]
\[
\tilde{\varphi}\bm{r}\times\FrameBasis = \varphi\bm{r}\times\FrameBasis - \bm{r}\times\FrameBasis \otimes \left(\FrameBasis \times \hat{\bm{r}}\right) \tilde{\bm{\SectionLocation}}
\]
Using this in \cref{eq:shear-center-constraint} furnishes the linear system:
\begin{equation}
    \SectionIntegral{\tilde{\varphi}\,\bm{r}\times\FrameBasis}
    =\left(\SectionIntegral{\bm{r}\times\FrameBasis\otimes\FrameBasis\times\bar{\bm{r}}} \right)\tilde{\bm{\SectionLocation}}
    \label{eq:derive-shear-center-1}
\end{equation}
where \(\bar{\bm{r}} \triangleq \bm{r} - \bar{\bm{\SectionLocation}}\).
Also recall that 
\[
\SectionIntegral{\varphi_{\SectionLocation} \bm{r}\times\FrameBasis} 
\equiv \SectionIntegral{\varphi_{\SectionLocation} (\bm{r}-\bar{\bm{\SectionLocation}})\times\FrameBasis} 
%\equiv \bm{0}
\]
Thus we use \(\bar{\bm{r}}\) for \(\bm{r}\) in \cref{eq:derive-shear-center-1}
\begin{equation}
    \SectionIntegral{\tilde{\varphi}\,\bar{\bm{r}}\times\FrameBasis}
    =\left(\SectionIntegral{\bar{\bm{r}}\times\FrameBasis\otimes\FrameBasis\times\bar{\bm{r}}} \right)\tilde{\bm{\SectionLocation}}
\end{equation}
Pulling out the constant \(\FrameBasis\) 
\begin{equation}
    \SectionIntegral{\tilde{\varphi}\,\bar{\bm{r}}\times\FrameBasis}
    =\FrameBasis^{\times}\left(\SectionIntegral{\bar{\bm{r}}\otimes\bar{\bm{r}}} \right)\FrameBasis^\times\tilde{\bm{\SectionLocation}}
\end{equation}


\subsection{Finite Elements}\label{finite-elements}

\cite{surana1979isoparametric}

\begin{quote}
\[
\begin{gathered}
\bm{K}_{a b}=\int \left(\frac{\partial N_{i}}{\partial x} \frac{\partial N_{b}}{\partial x}+\frac{\partial N_{i}}{\partial y} \frac{\partial N_{b}}{\partial y}\right) \operatorname{det}[J] \mathrm{d} \xi \mathrm{d} \eta \\
\bm{f}_{a}=\SectionIntegral{\left(y \frac{\partial N_{i}}{\partial x}-x \frac{\partial N_{i}}{\partial y}\right) \operatorname{det}[J]
} \\
\left\{\delta_{a}\right\}=\left\{\psi_{a}\right\}
\end{gathered}
\]
\end{quote}


\begin{equation}
    \bm{f}\cdot\bm{w} = \FrameBasis \cdot \SectionIntegral{(\bm{B}\bm{w} \times \FrameBasis_{\beta} \bm{N}\bm{x}^e_{\beta})}
\end{equation}

\begin{equation}
    \bm{f}\cdot\bm{w} = -\FrameBasis \cdot \FrameBasis_{\beta}\times \SectionIntegral{ (\bm{N}\bm{x}^e_{\beta})\, \bm{B}\bm{w}}
\end{equation}